\chapter{Bilangan sampai 100}\label{ch:g1:tens}

Membilang tumpukan yang besar satu per satu memakan waktu lama.
Inilah trik hebatnya: ikatlah benda-benda menjadi bungkus berisi
\emph{sepuluh}, lalu bilanglah bungkusnya. Beginilah bilangan sampai
$100$ dibentuk --- dan inilah sebabnya bilangan itu ditulis dengan
dua angka.

\section{Ikatan sepuluhan}

\begin{definition}[Puluhan dan satuan]\label{def:g1:tens:def}
Satu \emph{puluhan} adalah satu ikat berisi sepuluh satuan. Bilangan
$34$ berarti $3$ puluhan dan $4$ satuan:
\[
34 = 10 + 10 + 10 + 4 .
\]
Angka di kiri membilang ikatannya, angka di kanan membilang satuan
yang lepas.
\end{definition}

\begin{omfigure}
\begin{tikzpicture}[scale=0.75]
  % tiga ikat berisi sepuluh lidi, masing-masing diikat dengan pita membulat
  \foreach \b in {0,1,2} {
    \begin{scope}[xshift=\b*2.4cm]
    \foreach \i in {0,...,9} \draw[omDef, very thick]
      ({0.16*\i},0) -- ({0.16*\i},1.7);
    \draw[omExo, thick, rounded corners=3pt]
      (-0.22,0.62) rectangle (1.66,1.05);
    \end{scope}
  }
  % empat lidi lepas
  \foreach \i in {0,...,3} \draw[omThm, very thick]
    (7.7+0.4*\i,0) -- (7.7+0.4*\i,1.7);
  \node[right, font=\small] at (9.7,0.85)
    {$3$ ikat dan $4$ lidi: $34$};
\end{tikzpicture}

{\small Tiga puluh empat lidi: tiga ikat berisi sepuluh (masing-masing
diikat dengan pita) dan empat lidi lepas. Tidak perlu membilang satu
per satu sampai $34$!}
\end{omfigure}

\begin{example}[Membaca bilangannya]\label{ex:g1:tens:reading}
$20$ adalah dua puluhan (dua puluh), $30$ tiga puluhan (tiga
puluh)\,\dots\ sampai $90$ (sembilan puluh) dan $100$ (seratus:
sepuluh puluhan!). Di antaranya: $47$ dibaca ``empat puluh tujuh''
--- empat puluh untuk $4$ puluhannya, tujuh untuk satuannya.
\end{example}

\begin{omfigure}
\begin{tikzpicture}[scale=0.9]
  \draw (0,0) grid[xstep=2.2, ystep=0.8] (4.4,1.6);
  \node[font=\small] at (1.1,1.2) {puluhan};
  \node[font=\small] at (3.3,1.2) {satuan};
  \node[font=\large, omDef] at (1.1,0.4) {$3$};
  \node[font=\large, omThm] at (3.3,0.4) {$4$};
\end{tikzpicture}

{\small Tabel nilai tempat dari $34$: angka $3$ menempati tempat
puluhan, angka $4$ menempati tempat satuan. Tabel kecil ini bertambah
satu kolom setiap tahun!}
\end{omfigure}

\begin{example}[Angka sama, bilangan berbeda]\label{ex:g1:tens:place}
$27$ dan $72$ memakai angka yang sama tetapi sangat berbeda: $27$
adalah $2$ puluhan dan $7$ satuan; $72$ adalah $7$ puluhan dan $2$
satuan --- jauh lebih besar! \emph{Tempat} sebuah angka menentukan
apa yang dibilangnya.
\end{example}

\section{Papan seratus}

\begin{omfigure}
\begin{tikzpicture}[scale=0.52]
  \foreach \r in {0,...,9} {
    \foreach \c in {0,...,9} {
      \pgfmathtruncatemacro{\n}{\r*10+\c+1}
      \draw[gray!50] (\c,-\r) +(-0.5,-0.5) rectangle +(0.5,0.5);
      \node[font=\footnotesize] at (\c,-\r) {\n};
    }
  }
  \draw[omThm, very thick] (3,-4) +(-0.5,-0.5) rectangle +(0.5,0.5);
  \draw[omDef, very thick] (3,-5) +(-0.5,-0.5) rectangle +(0.5,0.5);
\end{tikzpicture}

{\small Papan seratus: setiap baris adalah satu puluhan. Satu langkah
ke kanan berarti $+1$; satu langkah \emph{ke bawah} berarti $+10$
(dari $44$ merah lurus ke bawah menuju $54$ biru).}
\end{omfigure}

\begin{method}[Berpindah di papan]\label{met:g1:tens:grid}
\begin{enumerate}
  \item $+1$ / $-1$: satu langkah ke kanan / ke kiri;
  \item $+10$ / $-10$: satu langkah ke bawah / ke atas --- hanya
        angka puluhannya yang berubah: $37 + 10 = 47$;
  \item membilang sepuluh-sepuluh: $10, 20, 30, \dots$ --- ikutilah
        kolom terakhir papan lurus ke bawah.
\end{enumerate}
\end{method}

\begin{example}[Membilang loncat 2 dan 5]\label{ex:g1:tens:skip}
Loncat $2$: $2, 4, 6, 8, 10, 12, \dots$ (bilangan genap). Loncat
$5$: $5, 10, 15, 20, 25, \dots$ --- semuanya berakhir dengan $5$ atau $0$
dan membentuk pola yang cantik di papan. Membilang loncat menyiapkan
perkalian (\cref{ch:g2:mult}).
\end{example}

\section{Membandingkan bilangan yang lebih besar}

\begin{method}[Membandingkan lewat puluhan]\label{met:g1:tens:compare}
Puluhan yang lebih banyak menang: $61 > 58$ karena $6$ puluhan
mengalahkan $5$ puluhan --- walaupun $8$ lebih banyak daripada $1$!
Hanya ketika puluhannya sama, satuanlah yang menentukan: $43 < 47$.
\end{method}

\begin{example}\label{ex:g1:tens:compare}
Urutkan $35$, $53$, $29$: lihat puluhannya dulu ($2$, $3$, $5$):
\[
29 < 35 < 53 .
\]
\end{example}

\section{Latihan}

\begin{exercise}[$\star$]\label{exo:g1:tens:1}
Berapa puluhan dan satuan pada: $26$; \ $40$; \ $73$; \ $99$?
\end{exercise}

\begin{exercise}[$\star$]\label{exo:g1:tens:2}
Tulislah bilangannya: $5$ puluhan dan $2$ satuan; \ $8$ puluhan; \
$1$ puluhan dan $9$ satuan; \ $10$ puluhan.
\end{exercise}

\begin{exercise}[$\star$]\label{exo:g1:tens:3}
Tulislah dengan angka: enam puluh tiga; \ empat puluh; \ delapan
puluh lima; \ seratus.
\end{exercise}

\begin{exercise}[$\star$]\label{exo:g1:tens:4}
Hitunglah dengan berpindah di papan seratus: $23 + 10$; \quad
$67 + 10$; \quad $45 - 10$; \quad $80 - 10$.
\end{exercise}

\begin{exercise}[$\star$]\label{exo:g1:tens:5}
Lanjutkan: $10, 20, 30, ?, ?, ?$. \quad
$5, 10, 15, ?, ?, ?$. \quad $12, 14, 16, ?, ?, ?$.
\end{exercise}

\begin{exercise}[$\star$]\label{exo:g1:tens:6}
Salin dan lengkapi dengan $<$ atau $>$:
\[
34 \;?\; 43, \qquad
58 \;?\; 61, \qquad
70 \;?\; 69, \qquad
99 \;?\; 100 .
\]
\end{exercise}

\begin{exercise}[$\star$]\label{exo:g1:tens:7}
Urutkan dari yang terkecil ke yang terbesar: $46$; \ $19$; \ $64$; \
$91$; \ $41$.
\end{exercise}

\begin{exercise}[$\star$]\label{exo:g1:tens:8}
Seorang tukang kebun mengikat $52$ bunga menjadi ikatan yang berisi
sepuluh. Berapa ikat yang penuh, dan berapa bunga yang lepas?
\end{exercise}

\begin{exercise}[$\star\star$]\label{exo:g1:tens:9}
Di papan seratus, mulailah dari $36$; melangkahlah dua langkah ke
bawah dan satu langkah ke kiri. Kamu mendarat di mana? Perhitungan
apa yang baru saja kamu lakukan?

\end{exercise}
